\documentclass[12pt]{article}
\usepackage[utf8]{inputenc}
\usepackage[T1]{fontenc}
\usepackage{hyperref}
\title{Kuo Language Resources for K-12 Curriculum Development}
\author{Compiled by Claude Sonnet 5.5 (claude-sonnet-5-5) for Andr\'as Kornai}
\date{}
\begin{document}
\maketitle
\section{Introduction}
Kuo (Koh; ISO 639-3 xuo) is an Mbum language of southern Chad and northern Cameroon. Speaker estimates vary widely between sources (about 3,000 to 21,000). Very few resources turned up in the search, and no dictionary, primer or school materials were found. Note that the French study on standardising ``Kwa'{}'' (Bakwa') concerns a different, Bamileke-group language of southwest Cameroon and is not a Kuo resource.
\section{Primary Educational Resources}
None found.
\section{Linguistic and Reference Materials}
\begin{itemize}
\item Wikipedia, Kuo language, which lists a consonant and vowel orthography and four tones: \url{https://en.wikipedia.org/wiki/Kuo_language}
\item Max Planck EVA numerals page for Kuo: \url{https://lingweb.eva.mpg.de/channumerals/Kuo.htm}
\end{itemize}
\section{Religious and Community Resources}
Joshua Project people-group page for Kuo in Cameroon: \url{https://www.joshuaproject.net/index.php/people_groups/12864/CM}. Language page: \url{https://www.joshuaproject.net/index.php/languages/xuo}. No scripture or community organisation was identified in the search results.
\section{Digital and Audio Resources}
None found.
\section{Curriculum Development Context}
The orthography listed on Wikipedia is a starting point but should be checked against local practice. A curriculum would need new materials developed with Kuo speakers in Chad and Cameroon, and the claim that Kuo is taught in some schools could not be verified from these results.
\end{document}